% ==============================================================================
% Fichier : 02_fondamentaux.tex
% Description : Définitions, Concepts et Bases de l'Audit et du Contrôle
% ==============================================================================

\chapter{Définitions, Concepts et Bases de l'Audit et du Contrôle}\label{chap:02_fondamentaux}

Ce chapitre établit le vocabulaire commun de l'auditeur. Une ambiguïté sur les termes "contrôle", "risque" ou "objectif" peut compromettre la compréhension de toute la mission.

\section{Audit et Contrôle : Distinction Fondamentale}
% ------------------------------------------------------------------------------

Il est fréquent de confondre ces deux notions, pourtant distinctes.

\subsection{Définitions}

\begin{boxdef}
\textbf{L'Audit :} Activité d'examen d'un \textbf{processus} ou d'une organisation par un tiers indépendant, visant à évaluer la conformité, l'efficacité et l'optimisation.
\end{boxdef}

\begin{boxdef}
\textbf{Le Contrôle :} Action de vérification ponctuelle d'une \textbf{opération} pour s'assurer de sa régularité, exactitude ou sincérité.
\end{boxdef}

\subsection{Comparaison}

\begin{table}[h!]
\centering
\begin{tabular}{p{3cm}p{5cm}p{5cm}}
    \toprule
    \textbf{Critère} & \textbf{Audit} & \textbf{Contrôle} \\
    \midrule
    \textbf{Sujet} & Un processus global. & Une opération unitaire. \\
    \textbf{But} & Amélioration, Optimisation. & Vérification, Conformité. \\
    \textbf{Acteur} & Indépendant (Interne ou Externe). & Souvent opérationnel. \\
    \textbf{Périodicité} & Ponctuelle (Mission). & Continue ou fréquente. \\
    \bottomrule
\end{tabular}
\caption{Différences entre Audit et Contrôle}
\end{table}

\begin{boxlizbiz}
\textbf{Exemple Lizbiz's :}
\begin{itemize}
    \item \textbf{Contrôle :} Le comptable vérifie que le montant d'une facture correspond à la commande reçue.
    \item \textbf{Audit :} L'auditeur analyse l'ensemble du processus "Achat-Fournisseur" pour voir si les contrôles (comme celui du comptable) sont efficaces et appliqués tout le temps.
\end{itemize}
\end{boxlizbiz}

% ------------------------------------------------------------------------------
\section{L'Assurance Raisonnable}
% ------------------------------------------------------------------------------

\begin{boxdef}
\textbf{Assurance Raisonnable :} Niveau de confiance élevé, mais non absolu, que l'auditeur apporte sur la fiabilité des éléments analysés.
\end{boxdef}

L'auditeur ne peut pas garantir à 100\% qu'il n'y a pas d'erreur ou de fraude (limite intrinsèque de l'échantillonnage, limites du temps, dissimulation possible). La certification ISO n'est pas une assurance "tous risques".

% ------------------------------------------------------------------------------
\section{Le Concept de Risque}
% ------------------------------------------------------------------------------

L'audit est intrinsèquement lié à la gestion des risques.

\subsection{Définition Académique}

Le risque est la probabilité qu'une menace exploite une vulnérabilité pour générer un impact sur une ressource.
\[ Risque = Probabilité \times Impact \]

\subsection{Le Cycle du Risque}

L'auditeur intervient souvent sur le "Risque Résiduel".

\begin{enumerate}
    \item \textbf{Risque Inhérent :} Risque brut, sans aucune mesure de protection.
    \item \textbf{Mesures de Contrôle :} Ce que l'entreprise a mis en place (Prévention, Détection).
    \item \textbf{Risque Résiduel :} Risque qui reste après application des contrôles. C'est celui que l'auditeur évalue.
\end{enumerate}

\begin{boxlizbiz}
\textbf{Cas Lizbiz's :}
Le \textbf{Risque Inhérent} de fraude à la commande est élevé. Lizbiz's a mis en place un \textbf{Contrôle} de validation manuelle par le manager. L'auditeur vérifie si ce contrôle fonctionne bien. Si le manager valide tout sans regarder, le \textbf{Risque Résiduel} reste élevé.
\end{boxlizbiz}

% ------------------------------------------------------------------------------
\section{Le Contrôle Interne (CI)}
% ------------------------------------------------------------------------------

\subsection{Définition}

\begin{boxdef}
\textbf{Contrôle Interne :} Ensemble des mesures prises par une organisation pour maîtriser ses activités et atteindre ses objectifs (Performance, Sécurité, Conformité).
\end{boxdef}

\subsection{Les Types de Contrôles}

On distingue généralement trois catégories de contrôles :
\begin{enumerate}
    \item \textbf{Préventifs :} Visent à empêcher l'erreur ou la fraude (ex: Mot de passe, Séparation des tâches).
    \item \textbf{Détectifs :} Visent à identifier l'erreur lorsqu'elle s'est produite (ex: Relevé de banque, Logs, Alarmes).
    \item \textbf{Correctifs :} Visent à corriger l'erreur et rétablir la situation (ex: Procédure de rejet de facture, Restauration de sauvegarde).
\end{enumerate}

% ------------------------------------------------------------------------------
\section{Principes d'Organisation du Contrôle}
% ------------------------------------------------------------------------------

Deux principes classiques régissent l'organisation du contrôle interne.

\subsection{La Règle des Quatre Yeux (Four Eyes Principle)}

\begin{boxdef}
Principe selon lequel toute décision ou opération sensible doit être validée par \textbf{deux personnes distinctes} et compétentes.
\end{boxdef}

Cela empêche une seule personne de détourner des fonds ou de valider ses propres erreurs.

\subsection{L'Allégorie des Deux Mains}

Pour empêcher la fraude, il faut séparer les tâches conflictuelles :
\begin{itemize}
    \item \textbf{Main Droite (Opérationnel) :} Celui qui exécute l'opération (ex: saisir la commande).
    \item \textbf{Main Gauche (Contrôle) :} Celui qui valide ou contrôle l'opération (ex: valider le paiement).
\end{itemize}

\begin{boxalert}
\textbf{Danger :} Si une même personne a le rôle "Créer Fournisseur" ET "Payer Fournisseur", le système est vulnérable à la fraude (création d'un fournisseur fictif). C'est une Non-Conformité majeure.
\end{boxalert}

% ------------------------------------------------------------------------------
\section{La Boîte à Outils de l'Auditeur : Le Vocabulaire de la Mission}
% ------------------------------------------------------------------------------

Pour structurer ses travaux, l'auditeur utilise un vocabulaire précis pour décrire ce qu'il fait.

\begin{table}[h!]
\centering
\small
\begin{tabular}{p{4cm}p{8cm}}
    \toprule
    \textbf{Terme} & \textbf{Définition} \\
    \midrule
    \textbf{Point de Contrôle} & Élément précis sur lequel porte l'audit (ex: Le serveur de paiement). \\
    \textbf{Objectif de Contrôle} & But à atteindre (ex: Assurer la confidentialité des données). \\
    \textbf{Critère d'évaluation} & La règle de référence (ex: La politique de sécurité, la norme ISO). \\
    \textbf{Question d'évaluation} & La question technique posée pour vérifier le critère (ex: "Les flux sont-ils chiffrés ?"). \\
    \textbf{Constatation} & Le résultat factuel de l'audit (ex: "Les flux ne sont pas chiffrés"). \\
    \bottomrule
\end{tabular}
\caption{Vocabulaire de la mission d'audit}
\end{table}

\subsection{Le "Contradictoire"}

Principe fondamental de l'audit : l'audité doit pouvoir s'expliquer et contester les faits avant la rédaction finale du rapport. L'auditeur ne peut pas porter de jugement sans avoir entendu l'audité.

\subsection{Observation vs Constatation}

\begin{itemize}
    \item \textbf{Constatation :} Un fait brut noté lors de l'enquête (ex: "L'employé n'a pas verrouillé sa session").
    \item \textbf{Observation :} L'observation validée et documentée dans le rapport, avec analyse de l'impact.
\end{itemize}

\subsection{Sur place et Sur pièces}

Deux méthodes de vérification :
\begin{itemize}
    \item \textbf{Sur place :} Présence physique de l'auditeur (entretiens, observation des locaux).
    \item \textbf{Sur pièces :} Analyse de documents à distance (revue de procédures, extraction de données).
\end{itemize}

% ------------------------------------------------------------------------------
\section{Définitions des Systèmes d'Information (SI)}
% ------------------------------------------------------------------------------

Pour auditer l'informatique, il faut en comprendre les composantes.

\subsection{Système Informatique vs Système d'Information}

\begin{boxdef}
\begin{itemize}
    \item \textbf{Système d'Information (SI) :} Ensemble organisé de ressources (personnes, données, procédures) permettant de traiter l'information dans l'organisation.
    \item \textbf{Système Informatique :} Partie technique du SI (Matériel, Logiciel, Réseau).
\end{itemize}
\end{boxdef}

L'auditeur audite le SI (le métier, l'usage) en s'appuyant sur la technique du Système Informatique.

\subsection{Composantes Techniques (Rappels)}

\begin{itemize}
    \item \textbf{Système d'exploitation :} Interface entre matériel et logiciels (Windows, Linux).
    \item \textbf{Base de données :} Structuration des données (PostgreSQL, Oracle).
    \item \textbf{Réseau :} Liaison entre les équipements.
    \item \textbf{Sécurité :} Protection des actifs (Firewall, Antivirus).
\end{itemize}
